% Part: first-order-logic
% Chapter: syntax-and-semantics
% Section: Structures

\documentclass[../../../include/open-logic-section]{subfiles}

\begin{document}

\olfileid{fol}{syn}{str}

\olsection{\printtoken{P}{structure} voor eerste-ordetalen}

\begin{explain}
Eerste-ordetalen zijn op zichzelf \emph{ongeïnterpreteerd}: aan de
!!{constant}s, !!{function}s en !!{predicate}s is geen specifieke
betekenis toegekend. Betekenis wordt toegekend door
\article{structure} \emph{!!{structure}} vast te leggen. Deze bepaalt
het \emph{domein}, dat wil zeggen de objecten waarnaar de !!{constant}s
verwijzen, waarop de !!{function}s werken en waarover de kwantoren
lopen. Daarnaast bepaalt zij welke !!{constant}s naar welke objecten
verwijzen, hoe !!a{function} objecten op objecten afbeeldt en op welke
objecten de !!{predicate}s van toepassing zijn. !!^{structure}s vormen
de grondslag voor \emph{semantische} begrippen in de logica, zoals
semantisch gevolg, geldigheid en vervulbaarheid. In de literatuur
worden zij afwisselend ``structuren,'' ``interpretaties'' of
``modellen'' genoemd.
\end{explain}

\begin{defn}[!!^{structure}s]
\Article{structure} \emph{!!{structure}}~$\Struct M$ voor een
eerste-ordetaal $\Lang{L}$ bestaat uit de volgende onderdelen:
\begin{enumerate}
\item \emph{Domein:} een niet-lege verzameling, $\Domain M$
\item \emph{Interpretatie van !!{constant}s:} voor iedere
  !!{constant}~$c$ van $\Lang{L}$ !!a{element}
  $\Assign{c}{M} \in \Domain M$
\item \emph{Interpretatie van !!{predicate}s:} voor ieder $n$-plaatsig
  !!{predicate}~$R$ van $\Lang{L}$ (behalve $\eq$) een $n$-plaatsige
  relatie $\Assign{R}{M} \subseteq \Domain{M}^n$
\item \emph{Interpretatie van !!{function}s:} voor ieder $n$-plaatsig
  !!{function}~$f$ van $\Lang{L}$ een $n$-plaatsige functie
  $\Assign{f}{M} \colon \Domain{M}^n \to \Domain{M}$
\end{enumerate}
\end{defn}

\begin{ex}
!!^a{structure}~$\Struct M$ voor de taal van de rekenkunde bestaat uit
  een verzameling, een element van $\Domain M$, $\Assign{\Obj 0}{M}$,
  als interpretatie van de !!{constant}~$\Obj 0$, een eenplaatsige functie
  $\Assign{\Obj \prime}{M} \colon \Domain{M} \to \Domain M$, twee
  tweeplaatsige functies $\Assign{\Obj +}{M}$ en $\Assign{\Obj
    \times}{M}$, beide $\Domain M^2 \to \Domain M$, en een
  tweeplaatsige relatie $\Assign{\Obj <}{M} \subseteq \Domain{M}^2$.

Een voor de hand liggend voorbeeld van zo'n structuur is:
\begin{enumerate}
\item $\Domain N = \Nat$
\item $\Assign{\Obj 0}{N} = 0$
\item $\Assign{\Obj \prime}{N}(n) = n + 1$ voor alle $n \in \Nat$
\item $\Assign{\Obj +}{N}(n, m) = n + m$ voor alle $n, m \in \Nat$
\item $\Assign{\Obj \times}{N}(n, m) = n\cdot m$ voor alle $n, m \in \Nat$
\item $\Assign{\Obj <}{N} = \Setabs{\tuple{n, m}}{n \in \Nat, m \in
  \Nat, n < m}$
\end{enumerate}
De zo gedefinieerde structuur~$\Struct N$ voor $\Lang L_A$ heet het
\emph{standaardmodel van de rekenkunde}, omdat zij de niet-logische
constanten van~$\Lang L_A$ precies interpreteert zoals men verwacht.

Er zijn echter nog veel andere !!{structure}s voor~$\Lang L_A$
mogelijk. We zouden bijvoorbeeld de verzameling~$\Int$ van gehele
getallen in plaats van~$\Nat$ als domein kunnen nemen en de
interpretaties van $\Obj 0$, $\Obj \prime$, $\Obj +$, $\Obj \times$ en
$\Obj <$ dienovereenkomstig kunnen definiëren. We kunnen voor~$\Lang
L_A$ evenwel ook structuren definiëren die zelfs in de verste verte
niets met getallen te maken hebben.
\end{ex}

\begin{ex}
Voor een structuur~$\Struct M$ voor de verzamelingenleerstaal~$\Lang
L_Z$ zijn slechts een verzameling en één tweeplaatsige relatie nodig.
Technisch gezien zou bijvoorbeeld de verzameling mensen samen met de
relatie ``$x$ is ouder dan $y$'' als !!a{structure} voor $\Lang L_Z$
kunnen dienen, evenals $\Nat$ samen met $n \ge m$ voor $n, m \in \Nat$.

Een bijzonder interessante !!{structure} voor $\Lang L_Z$, waarin de
!!{element}s van het domein zelf verzamelingen zijn en de interpretatie
van $\Obj \in$ daadwerkelijk de relatie ``$x$ is !!a{element} van~$y$''
is, is de !!{structure}~$\Struct{{HF}}$ van \emph{erfelijk eindige
verzamelingen}:
\begin{enumerate}
\item $\Domain{{{HF}}} = \emptyset \cup \Pow{\emptyset} \cup
  \Pow{\Pow{\emptyset}} \cup \Pow{\Pow{\Pow{\emptyset}}} \cup \dots$;
\item $\Assign{\Obj \in}{{{HF}}} = \Setabs{\tuple{x, y}}{x, y \in
  \Domain{{{HF}}}, x \in y}$.
\end{enumerate}
\end{ex}

\begin{digress}
De stipulaties die bepalen wat als !!a{structure} geldt, zijn van
invloed op onze logica. De uitsluiting van lege domeinen garandeert bij
de gebruikelijke definitie van vervulling (of waarheid) voor
gekwantificeerde zinnen bijvoorbeeld dat
$\lexists[x][(!A(x) \lor \lnot !A(x))]$ geldig is---dus een logische
waarheid. De stipulatie dat alle !!{constant}s naar een object in het
domein moeten verwijzen, garandeert bovendien dat existentiële
generalisatie een correct gevolgtrekkingspatroon is: $!A(a)$, dus
$\lexists[x][!A(x)]$. Als we zouden toestaan dat namen naar iets buiten
het domein verwijzen of in het geheel niet verwijzen, komen we in de
richting van een \emph{vrije logica}. Daar vereist existentiële
generalisatie een aanvullende premisse: $!A(a)$ en
$\lexists[x][\eq[x][a]]$, dus $\lexists[x][!A(x)]$.
\end{digress}

\end{document}
